\section{Related Work}
\label{sec:related}
Recent advancements in the post-training of large language models have been significantly driven by reinforcement learning with verifiable rewards (RLVR). This section systematically reviews several key areas that are highly relevant to and interconnected with the present study.
\subsection{Policy Optimization in RLVR.} 
Following DeepSeek-R1's validation of critic-free reinforcement learning \citep{guo2025deepseek}, Group Relative Policy Optimization (GRPO)\citep{shao2024deepseekmath} has emerged as a compute-efficient alternative to PPO \citep{schulman2017proximal}. However, its instability under long contexts and sparse rewards has spurred various enhancements. To mitigate policy staleness and granularity mismatch, CISPO clips importance sampling weights \citep{minimax2025minimaxm1scalingtesttimecompute} and GSPO employs sequence-level ratio calculations \citep{gspo}. For long-horizon and agentic tasks, DAPO introduces dynamic sampling and decoupled clipping \citep{Yu2025a} and  SAPO utilizes a temperature-controlled smooth gating function to form a continuous trust region \citep{SAPO2025}. Orthogonally, our work addresses GRPO's fundamental limitation of reward signal exhaustion by decoupling advantage normalization at the objective level. Since \method{} modifies the advantage computation rather than the optimization procedure, it can be composed with these enhancements, as we demonstrate with DAPO in \S\ref{sec:main-results}.


\subsection{Rubric-Based Training.}
Traditional outcome reward models frequently suffer from reward hacking, exploiting miracle steps to achieve correct final answers via flawed logic \citep{yuan2025miraclesteps}. To enforce rigorous logical constraints without prohibitive step-level human annotation, rubric-based process supervision has gained traction. Generative rubric reward models explicitly penalize disconnected derivations, drastically curtailing lucky guesses \citep{liang-etal-2025-generative}. In pursuit of self-verification, DeepSeekMath-V2 constructed a specialized LLM verifier to reversely drive the generator to patch reasoning loopholes \citep{deepseek-math-v2}. Furthermore, in complex open-domain tasks, frameworks like Agent-RRM propose multi-faceted evaluation systems enforcing explicit reasoning traces and actionable critiques \citep{fan2026exploring}. These studies establish that effective process supervision necessitates multi-dimensional criteria and explicit negative penalties. Building upon these insights, our work leverages a rubric-based process evaluation but introduces a correct-subset normalization mechanism to prevent verbosity-driven reward hacking.

\subsection{Generative Reward Models.}


To overcome the expressivity limitations of traditional discriminative models, reward modeling is shifting toward a generative paradigm \citep{Mahan2024}. By autoregressively generating deep comparative analyses prior to preference scoring, generative reward models significantly enhance interpretability and accuracy. This paradigm excels in abstract domains like formal proofs, as demonstrated by Proof-RM \citep{Yang2026}, and enables verifier-free reinforcement learning extrapolation in broad domains, as seen in RLPR \citep{Yu2025b}. However, generative evaluators remain susceptible to intrinsic biases such as imperfect sensitivity and specificity. To ensure statistically rigorous training, recent frameworks incorporate bias-correction estimators and adaptive calibration to construct statistically sound confidence intervals \citep{Lee2025}. Aligning with this generative shift, our methodology utilizes a generative judge to assess reasoning quality, yet uniquely confines its influence within the advantage space of correct solutions to immunize the policy against structural biases.


















% \paragraph{GRPO and policy optimization for LLMs.}
% GRPO \citepp{shao2024deepseekmath} simplifies RL-based LLM training by eliminating the critic network, computing advantages through group-level normalization of rewards from multiple responses to the same prompt. This approach was adopted by DeepSeek-R1 \citepp{guo2025deepseek} for training reasoning models at scale. DAPO \citepp{yu2025dapo} extends GRPO with techniques such as dynamic sampling and overlong penalty to improve training stability. 
% Scaf-GRPO \citepp{zhang2026scafgrpo} addresses the "learning cliff" in complex tasks by dynamically injecting hints to restore non-zero advantage gradients. Additionally, GMPO \citepp{zhao2025geometric} improves update stability by maximizing the geometric mean of token-level rewards rather than the arithmetic mean to suppress outliers. 
% Our work is complementary to these efforts: \method{} modifies the advantage computation rather than the optimization procedure, and can be combined with other GRPO enhancements.

% \paragraph{Process reward models.}
% \citept{lightman2023lets} introduced step-level process supervision for mathematical reasoning, demonstrating that verifying individual reasoning steps improves over outcome-only supervision. Math-Shepherd \citepp{wang2024math} automated process reward annotation using model-generated data. 
% OmegaPRM \citepp{ma2024omegaprm} fully automates step-level error isolation using Monte Carlo Tree Search. Concurrently, Process Advantage Verifiers \citepp{setlur2024rewardingprogressscalingautomated} reframe process rewards to measure probabilistic progress. 
% These works focus on \emph{learned} PRMs trained on step-level labels. In contrast, we use \emph{rubric-based} PRMs via LLM-as-Judge, which require no step-level annotations but introduce the reward hacking vulnerability that \method{} addresses. Our analysis of PRM reward hacking applies specifically to this rubric-based setting, where verbosity can inflate quality assessments.

% \paragraph{Reward shaping and signal design.}
% The challenge of designing informative reward signals for RL has a long history \citepp{ng1999policy}. In the LLM context, reward hacking, where models exploit reward model weaknesses to obtain high rewards without genuine improvement, has been documented as a persistent challenge \citepp{skalse2022defining}. While recent efforts attempt to mitigate this using modular Rubrics as Rewards (RaR) \citepp{gunjal2026rubrics}  and adaptive length penalties to combat verbosity bias \citepp{bu-etal-2025-beyond} , our work contributes a specific instance of this phenomenon (PRM reward hacking in GRPO) and a structural solution (decoupled normalization) rather than reward model improvement.

% \paragraph{LLM-as-Judge.}
% Using capable LLMs to evaluate other models' outputs has become a standard evaluation paradigm \citepp{zheng2024judging}, evolving towards fully self-rewarding architectures \citepp{yuan2024selfrewarding}.  We extend this paradigm to RL training by using LLM-based rubric evaluation as a process reward signal. 
% %Our finding that direct use of LLM-as-Judge scores for RL training causes reward hacking is relevant to the broader adoption of this paradigm: the same quality assessments that work well for evaluation can be pathological when used as RL rewards without careful integration.
% Our finding that direct use of LLM-as-Judge scores for RL training causes reward hacking is relevant to the broader adoption of this paradigm: the same quality assessments that work well for evaluation can be pathological when used as RL rewards without careful integration.


% \paragraph{Actor-critic methods.}
% In classical RL, actor-critic methods \citepp{schulman2017proximal} use a learned value function (critic) to reduce variance in policy gradient estimation. GRPO deliberately removes this critic, relying on group normalization instead. \method{} reintroduces a critic function, not a learned value network, but a rubric-based evaluator that assesses reasoning quality. This connects to the broader insight that some form of value assessment (beyond binary outcome) is needed for effective policy optimization, particularly as the policy improves and outcome signals become sparse.
